\section{估值结论：证据更强，赔率反而更窄}

新增合同和股东会通过提高了成功概率；购买资产发股、配套融资、官方东数一号权益价值和高杠杆降低了每股价值。旧43元目标撤销。

\begin{exhibitbox}[估值桥]
\noindent\begin{tabularx}{\textwidth}{L{3.4cm}R{1.8cm}R{1.8cm}X}
\toprule
\textbf{锚} & \textbf{数值} & \textbf{权重} & \textbf{说明} \\
\midrule
基本面锚 & 31.26元 & 55\% & 情景PE与官方交易SOTP混合，计入稀释 \\
Street锚 & 48.49元 & 25\% & 国金原始PDF 50.90元+国投38.86元 \\
市场事件锚 & 45.16元 & 20\% & 申万目标市值隐含45.16元 \\
最终目标 & 38.35元 & 100\% & 相对现价-1.6\% \\
\bottomrule
\end{tabularx}
\end{exhibitbox}

\section{完整估值公式}

本报告先计算三情景基本面价值，再用官方交易SOTP交叉验证，最后才引入Street和市场事件锚。订单金额不会直接加到市值：

\begin{exhibitbox}[五步估值复算]
\small
\noindent\begin{tabularx}{\textwidth}{L{3.4cm}X R{1.5cm}}
\toprule
\textbf{步骤} & \textbf{复算公式} & \textbf{结果} \\
\midrule
三情景概率值 & $22.10\times20\%+32.95\times55\%+48.87\times25\%$ & 34.76元 \\
官方交易SOTP & $(19.47\times40+115.00+0.80\times20)\div36.7509$ & 24.76元 \\
基本面锚 & $34.76\times65\%+24.76\times35\%$ & 31.26元 \\
Street锚 & $(50.90\times20\%+38.86\times5\%)\div25\%$ & 48.49元 \\
最终目标 & $31.26\times55\%+48.49\times25\%+45.16\times20\%$ & 38.35元 \\
\bottomrule
\end{tabularx}
\end{exhibitbox}

官方交易SOTP中的115亿元不是秦淮数据290.63亿元的100\%评估值，而是上市公司取得控制权所对应的东数一号全部权益交易锚；36.7509亿股为购买资产发股并假设80亿元配套融资按现价80\%发行后的基准股本。该公式同时约束资产价值和每股稀释。

\section{基准情景盈利桥}

\begin{exhibitbox}[2026E基准净利润与EPS]
\noindent\begin{tabularx}{\textwidth}{L{3.8cm}R{2.0cm}X}
\toprule
\textbf{组成} & \textbf{净利润信用} & \textbf{依据} \\
\midrule
传统材料／制冷剂／液冷 & 19.47亿元 & 国金2026-04-10原始PDF，未并表秦淮 \\
秦淮数据增量代理 & 6.64亿元 & 2025交易备考归母利润相对上市公司实际增量 \\
三笔算力合同信用 & 0.80亿元 & 年化未税收入中值80.19亿元，确认25\%，净利率4\% \\
House基准净利润 & 26.91亿元 & 上述三项合计 \\
基准稀释后股本 & 36.75亿股 & 购买资产发股+80亿元配套融资假设 \\
基准EPS & 0.732元 & 26.91亿元／36.75亿股 \\
\bottomrule
\end{tabularx}
\end{exhibitbox}

基准情景采用45倍PE，得到32.95元情景价值。该值不是最终目标，因为市场仍给予监管推进和AI平台化事件溢价，故再通过多锚模型得到38.35元。

\section{旧43元目标为什么下修}

\begin{exhibitbox}[目标价调整归因]
\noindent\begin{tabularx}{\textwidth}{L{3.2cm}L{2.2cm}X}
\toprule
\textbf{变化} & \textbf{方向} & \textbf{目标价影响} \\
\midrule
新增C单130--150亿元 & 上调 & 订单规模与客户验收证据增强 \\
A/B单已验收并确认收入 & 上调 & 算力信用由纯期权升级为小额基准盈利信用 \\
主业净利分母17.69升至19.47亿元 & 上调 & 使用国金最新原始PDF预测 \\
购买资产新增4.09亿股 & 下调 & 并表利润必须由更大股本分享 \\
配套融资股本敏感度 & 下调 & 基准股本由30.10亿股升至36.75亿股 \\
秦淮期权改用115亿元交易权益锚 & 下调 & 撤销旧模型任意510亿元期权价值 \\
72.39\%备考负债率、163.14亿元商誉 & 下调 & 提高财务和减值折价 \\
\bottomrule
\end{tabularx}
\end{exhibitbox}

综合结果是：经营证据变强，但每股价值和安全边际没有同步变强。旧43元目标因此撤销，新目标为38.35元。

\section{Street比较}

国金3月15日原始PDF目标50.90元、2026E EPS 0.636元、80x PE；其4月10日原始PDF把未并表秦淮的净利上修至19.47亿元、EPS 0.647元，但未给新目标。国投38.86元与富途均值39.68元更接近现价。House最终目标38.35元低于旧43元，原因是官方交易权益和稀释替代了任意期权估值。

\section{操作策略}

已有仓位：持有但不追加叙事溢价；监管推进、配套融资定价和C单验收是加仓门槛。新开仓：等待37.15元事件缺口或34.77--35元深度支撑。45.16元以上只在硬公告与成交共振时做事件交易。

\sourcenote{data/current\_valuation\_model\_20260713.json；data/transaction\_dilution\_model\_20260713.json；两份国金原始PDF。}
